\documentclass[12pt]{article}
\usepackage[utf8]{inputenc}
\usepackage[T1]{fontenc}
\usepackage[margin=1in]{geometry}
\usepackage{setspace}
\usepackage{xcolor}
\usepackage[colorlinks=true, urlcolor=blue!60!black, linkcolor=blue!60!black]{hyperref}
\setlength{\parindent}{2em}
\setlength{\parskip}{1.25em}
\renewcommand{\baselinestretch}{1.0}
\begin{document}
\begin{center}
{\LARGE \textit{United States v. Rahimi}, 602 U.S. 680 (2024)\par}
\vspace{0.8em}
{\large Case Brief\par}
\end{center}
\vspace{1em}

\section{Memory Jogger}

In an 8--1 decision clarifying \textit{Bruen}, the Supreme Court upheld 18 U.S.C.\ \S\ 922(g)(8) against a facial Second Amendment challenge, holding that an individual found by a court to pose a credible threat to the physical safety of another may be temporarily disarmed consistent with the principles underpinning the Nation's historical tradition of firearm regulation (\href{https://www.courtlistener.com/opinion/10600058/united-states-v-rahimi/}{CourtListener}; \href{https://www.law.cornell.edu/supremecourt/text/22-915}{Cornell LII}).

\section{Detailed Case Facts}

In December 2019, Zackey Rahimi met his girlfriend, C. M., who was also the mother of his young child, A. R., for lunch in a parking lot in Texas (\href{https://www.law.cornell.edu/supremecourt/text/22-915}{Cornell LII}; \href{https://www.supremecourt.gov/opinions/23pdf/22-915_8o6b.pdf}{Slip Op.\ at 1}). When an argument erupted and C. M.\ tried to leave, Rahimi grabbed her by the wrist, dragged her back to his car, and shoved her inside, causing her head to strike the dashboard. Noticing a bystander watching the assault, Rahimi retrieved a firearm from under the passenger seat, allowing C. M.\ to flee; Rahimi fired his gun as she ran away, and later called her to threaten that he would shoot her if she reported the incident (\href{https://www.law.cornell.edu/supremecourt/text/22-915}{Cornell LII}).

C. M.\ applied for a civil protective order, submitting an affidavit recounting the parking-lot assault, prior assaults, and conduct endangering their child. Although given notice and an opportunity to contest the allegations, Rahimi did not do so. On February 5, 2020, a Tarrant County, Texas state court issued a two-year protective order with both parties' consent, expressly finding that Rahimi had committed ``family violence,'' that such violence was ``likely to occur again,'' and that he posed ``a credible threat'' to the ``physical safety'' of C. M.\ or A. R.\ (\href{https://www.law.cornell.edu/supremecourt/text/22-915}{Cornell LII}). The order barred Rahimi from threatening C. M.\ or her family, restricted contact with C. M.\ to matters concerning their child, prohibited him from using or threatening physical force, and suspended his gun license for two years (\href{https://www.law.cornell.edu/supremecourt/text/22-915}{Cornell LII}).

Rahimi repeatedly violated the order. In May 2020, he approached C. M.'s home at night and contacted her through social media. In November 2020, he threatened a different woman with a gun, leading to an aggravated assault charge. Between December 2020 and January 2021, Arlington, Texas police linked him to five additional shootings, including firing into a drug customer's home, shooting at another motorist after a collision, firing into the air in a residential neighborhood, firing at a truck after a highway road-rage chase, and firing into the air at a roadside burger restaurant when his friend's credit card was declined (\href{https://www.law.cornell.edu/supremecourt/text/22-915}{Cornell LII}; \href{https://www.oyez.org/cases/2023/22-915}{Oyez}). Executing a search warrant at Rahimi's residence, police recovered a pistol, a rifle, ammunition, and a copy of the protective order (\href{https://www.law.cornell.edu/supremecourt/text/22-915}{Cornell LII}).

\section{Procedural History}

A federal grand jury indicted Rahimi in the U.S. District Court for the Northern District of Texas on one count of possessing a firearm while subject to a domestic violence restraining order in violation of 18 U.S.C.\ \S\ 922(g)(8)---then punishable by up to 10 years' imprisonment under \S\ 924(a)(2), later increased to 15 years under \S\ 924(a)(8) by the Bipartisan Safer Communities Act, Pub.\ L.\ No.\ 117-159 (\href{https://www.law.cornell.edu/supremecourt/text/22-915}{Cornell LII}). Section 922(g)(8) applies only when (A) the order was issued after a hearing with actual notice and an opportunity to participate; (B) the order restrains the defendant from harassing, stalking, or threatening an intimate partner or child; and (C) the order either (i) includes a finding that the defendant ``represents a credible threat to the physical safety'' of the intimate partner or child, or (ii) explicitly prohibits the use, attempted use, or threatened use of physical force (\href{https://www.law.cornell.edu/supremecourt/text/22-915}{Cornell LII}).

Rahimi moved to dismiss the indictment, bringing a facial Second Amendment challenge to \S\ 922(g)(8). The District Court denied the motion as foreclosed by Fifth Circuit precedent (\textit{United States v. McGinnis}, 956 F.3d 747 (5th Cir.\ 2020)), and Rahimi pleaded guilty while preserving his constitutional challenge (\href{https://harvardlawreview.org/print/vol-138/united-states-v-rahimi/}{138 Harv.\ L.\ Rev.\ 325}). On June 8, 2022, a Fifth Circuit panel affirmed in an unpublished \textit{per curiam} opinion (\href{https://harvardlawreview.org/print/vol-138/united-states-v-rahimi/}{138 Harv.\ L.\ Rev.\ 325}). While Rahimi's petition for rehearing en banc was pending, the Supreme Court decided \textit{New York State Rifle \& Pistol Ass'n v. Bruen}, 597 U.S. 1 (2022), repudiating means-end scrutiny in favor of a text-and-historical-tradition test. The Fifth Circuit withdrew its opinion, ordered supplemental briefing, and reversed, vacating Rahimi's conviction on the ground that the Government had not identified a sufficiently similar historical analogue for \S\ 922(g)(8) (61 F.4th 443 (5th Cir.\ 2023)) (\href{https://www.law.cornell.edu/supremecourt/text/22-915}{Cornell LII}). The Supreme Court granted certiorari on June 30, 2023, heard oral argument on November 7, 2023, and issued its decision on June 21, 2024 (\href{https://www.scotusblog.com/cases/united-states-v-rahimi/}{SCOTUSblog}).

\section{Judicial Votes}

\begin{itemize}
\item \textbf{Vote:} 8--1 (reversed and remanded).
\item \textbf{Majority Opinion:} Chief Justice John G. Roberts, Jr.\ delivered the opinion of the Court, joined in full by Justices Samuel A. Alito, Jr., Sonia Sotomayor, Elena Kagan, Neil M. Gorsuch, Brett M. Kavanaugh, Amy Coney Barrett, and Ketanji Brown Jackson (\href{https://www.law.cornell.edu/supremecourt/text/22-915}{Cornell LII}; \href{https://www.scotusblog.com/cases/united-states-v-rahimi/}{SCOTUSblog}).
\item \textbf{Concurring Opinions:}
  \begin{itemize}
  \item Justice Sonia Sotomayor filed a concurring opinion, joined by Justice Elena Kagan.
  \item Justice Neil M. Gorsuch filed a concurring opinion.
  \item Justice Brett M. Kavanaugh filed a concurring opinion.
  \item Justice Amy Coney Barrett filed a concurring opinion.
  \item Justice Ketanji Brown Jackson filed a concurring opinion.
  \end{itemize}
\item \textbf{Dissenting Opinion:} Justice Clarence Thomas filed a lone dissenting opinion (\href{https://www.law.cornell.edu/supremecourt/text/22-915}{Cornell LII}).
\end{itemize}

\section{Holding}

\textbf{Judgment:} The judgment of the U.S. Court of Appeals for the Fifth Circuit was \textbf{reversed}, and the case was \textbf{remanded} for further proceedings (\href{https://www.law.cornell.edu/supremecourt/text/22-915}{Cornell LII}).

\textbf{Rule Announced:}
\begin{enumerate}
\item \textit{Substantive Second Amendment Holding:} When an individual has been found by a court to pose a credible threat to the physical safety of another, that individual may be temporarily disarmed consistent with the Second Amendment. Consequently, 18 U.S.C.\ \S\ 922(g)(8)(C)(i) is constitutional as applied to Rahimi and therefore survives his facial challenge under \textit{United States v. Salerno}, 481 U.S. 739, 745 (1987) (\href{https://www.law.cornell.edu/supremecourt/text/22-915}{Cornell LII}).
\item \textit{Methodological Clarification of Bruen:} Under \textit{Bruen}, a modern firearm regulation need not be a ``dead ringer'' or ``historical twin'' of a 1791 statute; rather, courts ask whether the challenged law is consistent with the \textit{principles} that underpin the Nation's historical tradition of firearm regulation, examining ``why'' and ``how'' the law burdens the right (\href{https://www.law.cornell.edu/supremecourt/text/22-915}{Cornell LII}).
\item \textit{Rejection of the ``Responsible'' Citizen Test:} The Government may not disarm individuals simply by labeling them not ``responsible,'' a vague term never adopted as a legal boundary in \textit{Heller} or \textit{Bruen} (\href{https://www.law.cornell.edu/supremecourt/text/22-915}{Cornell LII}).
\end{enumerate}

\textbf{Holding vs.\ Dicta:} The Court's holding is expressly narrow: it decides only that \S\ 922(g)(8)(C)(i)---which requires an individualized judicial finding of a ``credible threat''---is constitutional as applied to Rahimi, defeating his facial challenge (\href{https://www.law.cornell.edu/supremecourt/text/22-915}{Cornell LII}). The Court expressly declined to decide: (1) the constitutionality of \S\ 922(g)(8)(C)(ii) (orders prohibiting physical force without an explicit credible-threat finding); (2) whether categorical legislative bans on classes of persons (such as nonviolent felons under \S\ 922(g)(1)) or permanent disarmament are constitutional, though it reiterated in dicta \textit{Heller}'s statement that prohibitions on felons and the mentally ill are ``presumptively lawful''; (3) whether 1791 or 1868 is the controlling historical anchor for Second Amendment interpretation (Footnote 1); and (4) any Fifth Amendment due process challenges to \S\ 922(g)(8) (Footnote 2) (\href{https://www.law.cornell.edu/supremecourt/text/22-915}{Cornell LII}; \href{https://harvardlawreview.org/print/vol-138/united-states-v-rahimi/}{138 Harv.\ L.\ Rev.\ 325}).

\section{Analysis of Opinions}

\subsection{Majority Opinion (Chief Justice Roberts)}

Chief Justice Roberts began by correcting what he termed lower courts' ``misunderst[anding]'' of \textit{District of Columbia v. Heller}, 554 U.S. 570 (2008), and \textit{Bruen}: the Court's Second Amendment precedents ``were not meant to suggest a law trapped in amber'' (\href{https://www.law.cornell.edu/supremecourt/text/22-915}{Cornell LII}). Just as the Second Amendment protects modern bearable arms beyond muskets and sabers, it permits modern firearm regulations that are ``relevantly similar'' to the \textit{principles} underpinning the Nation's historical regulatory tradition, evaluated by ``why'' and ``how'' they burden the right (\href{https://www.law.cornell.edu/supremecourt/text/22-915}{Cornell LII}).

Applying \textit{Salerno}'s demanding standard for facial challenges---under which the Government prevails if the statute is valid in at least some applications---Roberts focused exclusively on \S\ 922(g)(8)(C)(i) and identified two Founding-era regulatory traditions: \textbf{surety laws} and \textbf{``going armed'' (affray) laws} (\href{https://www.law.cornell.edu/supremecourt/text/22-915}{Cornell LII}). Surety laws, a form of preventive justice described by Blackstone and codified in Massachusetts and at least nine other jurisdictions, empowered magistrates after notice and a hearing to require persons suspected of future violence---including abusive spouses and those going armed offensively---to post a peace bond for a limited period or face jail. ``Going armed'' laws, descending from the Statute of Northampton (1328) and codified in at least four early states, punished individuals who terrified the public with dangerous weapons through forfeiture of arms and imprisonment (\href{https://www.law.cornell.edu/supremecourt/text/22-915}{Cornell LII}). Taken together, these regimes established a clear principle: an individual found by a court to pose a clear threat of physical violence to another may be temporarily disarmed. Section 922(g)(8)(C)(i) matched both the ``why'' (mitigating demonstrated threats of physical violence) and the ``how'' (temporary disarmament following an individualized judicial finding, a lesser burden than imprisonment) (\href{https://www.law.cornell.edu/supremecourt/text/22-915}{Cornell LII}). Roberts faulted the Fifth Circuit both for demanding a ``historical twin'' and for inverting \textit{Salerno} by chasing hypothetical unconstitutional applications, while also rejecting the Government's overbroad ``not responsible'' theory (\href{https://www.law.cornell.edu/supremecourt/text/22-915}{Cornell LII}).

\subsection{Concurring Opinion (Justice Sotomayor, joined by Justice Kagan)}

Justice Sotomayor, joined by Justice Kagan, joined the Court's opinion in full while reiterating that \textit{Bruen} was wrongly decided and that traditional means-end scrutiny remains the proper framework (\href{https://www.law.cornell.edu/supremecourt/text/22-915}{Cornell LII}). She praised the majority for clarifying that \textit{Bruen} looks to underlying historical \textit{principles} rather than exact matches, warning that Justice Thomas's rigid approach would turn \textit{Bruen} into a ``one-way ratchet'' that hamstrings modern democracy---especially on domestic violence, a problem Founding-era law largely ignored because women were excluded from the polity (\href{https://www.law.cornell.edu/supremecourt/text/22-915}{Cornell LII}). Citing empirical data that an abused woman is five times more likely to be murdered when an abuser has access to a gun, she concluded that \S\ 922(g)(8) easily survives under either \textit{Bruen} or heightened scrutiny (\href{https://www.law.cornell.edu/supremecourt/text/22-915}{Cornell LII}).

\subsection{Concurring Opinion (Justice Gorsuch)}

Justice Gorsuch emphasized the narrow scope of a facial challenge under \textit{Salerno} and defended originalist text-and-history methodology over pre-\textit{Bruen} interest-balancing, which he criticized as a ``black box regime'' enabling judicial policymaking (\href{https://www.law.cornell.edu/supremecourt/text/22-915}{Cornell LII}). Drawing an analogy to the Sixth Amendment Confrontation Clause in \textit{Crawford v. Washington}, 541 U.S. 36 (2004), he cautioned courts against extrapolating overly abstract ``values'' to water down enumerated rights. Crucially, Gorsuch underscored what the Court left open for future as-applied challenges: disarmament without a judicial credible-threat finding, permanent disarmament, enforcement against someone using a gun in self-defense, or categorical legislative bans on groups deemed ``not responsible'' (\href{https://www.law.cornell.edu/supremecourt/text/22-915}{Cornell LII}).

\subsection{Concurring Opinion (Justice Kavanaugh)}

Justice Kavanaugh wrote a sweeping methodological defense of interpreting vague constitutional text through \textbf{text, pre-ratification history, post-ratification history (tradition/liquidation), and precedent} rather than ``policy''-driven tiers of scrutiny (\href{https://www.law.cornell.edu/supremecourt/text/22-915}{Cornell LII}). Invoking James Madison (\textit{Federalist No.\ 37}), Chief Justice John Marshall (\textit{McCulloch v. Maryland}), and Justice Antonin Scalia, Kavanaugh argued that consistent post-ratification practice helps liquidate ambiguous provisions and constrains judges to act as neutral ``umpires'' rather than legislators (\href{https://www.law.cornell.edu/supremecourt/text/22-915}{Cornell LII}; \href{https://pennlawreview.com/2025/06/10/kavanaugh-rahimi-and-the-effort-to-legitimate-originalism/}{Penn Law Review}). He characterized tiered scrutiny as a modern, subjective invention of the 1950s and 1960s that should not be extended to the Second Amendment.

\subsection{Concurring Opinion (Justice Barrett)}

Justice Barrett focused on two core originalist questions. First, parting ways with Kavanaugh's broader embrace of ``tradition,'' she argued that under ``originalism \textit{simpliciter},'' only history illuminating original public meaning at ratification is binding; late post-enactment tradition unmoored from original meaning is not binding law (\href{https://www.law.cornell.edu/supremecourt/text/22-915}{Cornell LII}). Second, she addressed the ``level of generality'' problem in what she called ``original contours'' history: demanding a Founding-era twin forces 21st-century legislatures into 18th-century policy molds and wrongly assumes Founders exercised their full regulatory authority (``use it or lose it''). Instead, ``[h]istorical regulations reveal a principle, not a mold,'' and the majority chose the right level of generality here (\href{https://www.law.cornell.edu/supremecourt/text/22-915}{Cornell LII}).

\subsection{Concurring Opinion (Justice Jackson)}

Justice Jackson joined the opinion as a faithful application of binding precedent under \textit{Bruen}, while writing separately to sound an alarm about \textit{Bruen}'s unworkability (\href{https://www.law.cornell.edu/supremecourt/text/22-915}{Cornell LII}). Surveying a dozen lower-court decisions complaining of confusion and inconsistent outcomes, she observed that ``there is little method to \textit{Bruen}'s madness'' and that courts and legislators are being conscripted as amateur historians without answers to fundamental questions---such as which era controls, how many analogues constitute a tradition, and what level of generality applies---undermining the stability and predictability required by the Rule of Law (\href{https://www.law.cornell.edu/supremecourt/text/22-915}{Cornell LII}).

\subsection{Dissenting Opinion (Justice Thomas)}

Justice Thomas, the author of \textit{Bruen}, dissented alone, concluding that ``[n]ot a single historical regulation justifies'' \S\ 922(g)(8) (\href{https://www.law.cornell.edu/supremecourt/text/22-915}{Cornell LII}). Because Rahimi is among ``the people'' and \S\ 922(g)(8) bans firearm possession even in the home, the Government bore the burden of identifying a single historical regulation imposing a comparable burden for a comparable justification. Thomas argued that English ``dangerous person'' laws (such as the Militia Act of 1662) were designed to suppress political and religious rebellion---the very abuses that inspired the Second Amendment---and that failed Founding-era ratification proposals mentioning ``peaceable'' citizens carried no weight (\href{https://www.law.cornell.edu/supremecourt/text/22-915}{Cornell LII}). Turning to the majority's two analogues, Thomas contended that: (1) \textbf{surety laws} addressed interpersonal violence through a materially less severe means (a financial peace bond that left the individual's right to keep and bear arms intact, unlike \S\ 922(g)(8)'s total disarmament backed by a 15-year felony penalty); and (2) \textbf{affray (``going armed'') laws} were criminal prosecutions requiring proof beyond a reasonable doubt for public conduct terrifying the community, expressly excluding private domestic violence in the home (\href{https://www.law.cornell.edu/supremecourt/text/22-915}{Cornell LII}). Accusing the majority of improperly ``mixing and matching'' one law's justification with another law's burden, Thomas argued that States should disarm violent offenders through criminal prosecution and conviction, not civil restraining orders lacking Sixth Amendment safeguards (\href{https://www.law.cornell.edu/supremecourt/text/22-915}{Cornell LII}).

\section{Examples: Future Applications}

\begin{enumerate}
\item \textbf{Same-Side Example 1 (Extreme Risk Protection Order / ``Red Flag'' Law for Credible Threat to Others):} A state statute authorizes a judge, after notice and an adversarial hearing, to issue a one-year Extreme Risk Protection Order temporarily barring a respondent from possessing firearms upon finding by clear and convincing evidence that the respondent poses an imminent, credible threat of physical violence to coworkers or family members. \textit{Analysis:} Under \textit{Rahimi}, this law almost certainly survives a Second Amendment challenge because it mirrors \S\ 922(g)(8)(C)(i) in both ``why'' (preventing demonstrated physical violence against others) and ``how'' (temporary disarmament following an individualized judicial finding with procedural protections) (\href{https://jaapl.org/custom-print/68663}{JAAPL}).

\item \textbf{Same-Side Example 2 (State Restraining Order Disarming a Violent Stalker):} A state statute temporarily prohibits firearm possession by any individual subject to a civil anti-stalking protective order issued after a noticed hearing in which a court explicitly finds that the respondent represents a credible threat to the physical safety of a neighbor or coworker. \textit{Analysis:} Even though the victim is not an ``intimate partner,'' \textit{Rahimi}'s stated principle applies to any individual found by a court to pose ``a credible threat to the physical safety of another,'' and historical surety laws protected neighbors and strangers as well as spouses (\href{https://www.law.cornell.edu/supremecourt/text/22-915}{Cornell LII}).

\item \textbf{Opposite-Side Example 1 (Permanent Disarmament for a Nonviolent Regulatory Felony):} The federal government prosecutes a citizen under 18 U.S.C.\ \S\ 922(g)(1) for possessing a hunting rifle at home twenty years after a single nonviolent felony conviction for making a false statement on a food-stamp application, with no history or judicial finding of violence. \textit{Analysis:} This falls on the opposite side of \textit{Rahimi}'s reasoning because \S\ 922(g)(1) imposes \textit{permanent}, categorical disarmament without any individualized judicial finding of a credible threat of physical violence---the two features Roberts and Gorsuch repeatedly emphasized in upholding \S\ 922(g)(8)(C)(i) (\href{https://harvardlawreview.org/print/vol-138/united-states-v-rahimi/}{138 Harv.\ L.\ Rev.\ 325}; \href{https://virginialawreview.org/articles/united-states-v-rahimi-we-do-not-resolve-any-of-those-questions-because-we-cannot/}{Virginia Law Review}).

\item \textbf{Opposite-Side Example 2 (Legislative Disarmament of ``Irresponsible'' Citizens Based on Unpaid Debt or Traffic Infractions):} A state enacts a statute barring anyone who has accumulated three speeding tickets or defaulted on a civil judgment from possessing firearms for five years, arguing that such persons are not ``responsible'' citizens. \textit{Analysis:} \textit{Rahimi} squarely rejected the Government's ``not responsible'' theory as vague and unmoored from text and history (\href{https://www.law.cornell.edu/supremecourt/text/22-915}{Cornell LII}). Because civil debt default or routine traffic violations do not establish a credible threat of physical violence to another, the law lacks any historical analogue.

\item \textbf{Fence-Sitter Example (Automatic Disarmament Under \S\ 922(g)(8)(C)(ii) or Suicide-Only ERPOs):} Consider either (a) a federal prosecution under 18 U.S.C.\ \S\ 922(g)(8)(C)(ii) based on a boilerplate mutual divorce restraining order that bars ``physical force'' without any judicial finding that the respondent poses a credible threat, or (b) a civil Extreme Risk Protection Order temporarily disarming a person solely due to a risk of \textit{self-harm} (suicide) rather than violence against others. \textit{Analysis:} Both sit squarely on the fence after \textit{Rahimi}. The majority expressly reserved \S\ 922(g)(8)(C)(ii) (\href{https://www.law.cornell.edu/supremecourt/text/22-915}{Cornell LII}), and Gorsuch highlighted that disarmament without a credible-threat finding remains unresolved. Similarly, because \textit{Rahimi} grounded its historical principle in preventing physical harm ``to another,'' suicide-only red-flag orders test whether courts will accept a higher level of generality or require distinct Founding-era analogues (\href{https://jaapl.org/custom-print/68663}{JAAPL}).
\end{enumerate}

\section{Critique}

\textbf{Originalist and Textualist Critiques:} From the originalist right, scholars and commentators have noted that \textit{Rahimi} exposed deep methodological fractures among the Court's six conservative Justices---producing ``one, two \dots many originalisms'' (\href{https://houstonlawreview.org/article/147417-_united-states-v-rahimi_-five-takes}{Houston Law Review}). Strict originalists echo Justice Thomas's complaint that the majority quietly retreated from \textit{Bruen}---engaging in ``\textit{Bruen}-lite'' by shifting from concrete historical traditions to abstract ``principles'' and by stitching together the justification of surety laws with the penalty of affray laws (\href{https://harvardlawreview.org/print/vol-138/united-states-v-rahimi/}{138 Harv.\ L.\ Rev.\ 325}; \href{https://houstonlawreview.org/article/147417-_united-states-v-rahimi_-five-takes}{Houston Law Review}). Others observe that the clash between Kavanaugh's embrace of post-ratification tradition and Barrett's insistence on ratification-era ``originalism \textit{simpliciter}'' leaves unresolved whether ``tradition'' has independent constitutional authority or merely evidences original public meaning (\href{https://harvardlawreview.org/print/vol-138/united-states-v-rahimi/}{138 Harv.\ L.\ Rev.\ 325}; \href{https://www.floridalawreview.com/article/162313-tradition-and-discretion/attachment/345085.pdf}{Florida Law Review}).

\textbf{Progressive and Doctrinal Critiques:} Progressive and doctrinal scholars---including Joseph Blocher, Eric Ruben, Reva Siegel, and commentators in the \textit{Harvard}, \textit{Virginia}, and \textit{University of Pennsylvania Law Reviews}---argue that while \textit{Rahimi} avoided the catastrophic practical outcome of arming domestic abusers, it failed to fix \textit{Bruen}'s foundational flaws. First, shifting from ``historical tradition'' to ``principles that underpin our regulatory tradition'' fails to solve the \textbf{level-of-generality problem}: without a principled rule for how broadly or narrowly to abstract a historical principle, lower courts remain mired in subjective, ideologically correlated outcomes (\href{https://virginialawreview.org/articles/united-states-v-rahimi-we-do-not-resolve-any-of-those-questions-because-we-cannot/}{Virginia Law Review}; \href{https://scholarship.law.duke.edu/cgi/viewcontent.cgi?article=6973&context=faculty_scholarship}{Blocher \& Ruben, Duke Law Repository}). Second, feminist and equality scholars critique the Court's continued insistence on Founding-era analogues as structurally biased: tying modern gun safety to 18th-century lawmaking privileges an era when women and people of color were excluded from political participation and domestic abuse was tolerated under coverture (\href{https://repository.law.miami.edu/cgi/viewcontent.cgi?article=2286&context=fac_articles}{University of Miami Law Repository citing Siegel}). Third, critics contend that Kavanaugh's attack on tiered scrutiny masks the reality that selecting historical analogues and levels of generality involves inescapable value judgments (\href{https://pennlawreview.com/2025/06/10/kavanaugh-rahimi-and-the-effort-to-legitimate-originalism/}{Penn Law Review}).

\textbf{Underlying Competing Values:} \textit{Rahimi} reflects the deep constitutional collision between (1) \textit{individual armed self-defense and judicial constraint}---the concern that unmoored interest-balancing lets judges erode an enumerated right---and (2) \textit{democratic public-safety governance and equality}---the imperative that modern legislatures be able to protect vulnerable citizens, especially domestic-violence victims, from lethal modern firearms using contemporary legal tools.

\section{Key Quotations}

\begin{quote}
Nevertheless, some courts have misunderstood the methodology of our recent Second Amendment cases. These precedents were not meant to suggest a law trapped in amber.
\end{quote}
\noindent --- Chief Justice Roberts, opinion of the Court, 602 U.S.\ at 691 (slip op.\ at 7) (\href{https://www.law.cornell.edu/supremecourt/text/22-915}{Cornell LII}; \href{https://www.courtlistener.com/opinion/10600058/united-states-v-rahimi/}{CourtListener}).

\begin{quote}
Taken together, the surety and going armed laws confirm what common sense suggests: When an individual poses a clear threat of physical violence to another, the threatening individual may be disarmed.
\end{quote}
\noindent --- Chief Justice Roberts, opinion of the Court (slip op.\ at 13) (\href{https://www.law.cornell.edu/supremecourt/text/22-915}{Cornell LII}).

\begin{quote}
Rather, we conclude only this: An individual found by a court to pose a credible threat to the physical safety of another may be temporarily disarmed consistent with the Second Amendment.
\end{quote}
\noindent --- Chief Justice Roberts, opinion of the Court (slip op.\ at 17) (\href{https://www.law.cornell.edu/supremecourt/text/22-915}{Cornell LII}).

\begin{quote}
``Analogical reasoning'' under \textit{Bruen} demands a wider lens: Historical regulations reveal a principle, not a mold.
\end{quote}
\noindent --- Justice Barrett, concurring (slip op.\ at 4) (\href{https://www.law.cornell.edu/supremecourt/text/22-915}{Cornell LII}).

\begin{quote}
Not a single historical regulation justifies the statute at issue, 18 U.S.C.\ \S\ 922(g)(8). Therefore, I respectfully dissent.
\end{quote}
\noindent --- Justice Thomas, dissenting (slip op.\ at 1) (\href{https://www.law.cornell.edu/supremecourt/text/22-915}{Cornell LII}).


\end{document}
